\documentclass[a5paper,10pt]{book}

\usepackage{amsfonts}
\usepackage{hyperref}
\usepackage{bookmark}
% \usepackage[utf8]{inputenc}
\usepackage[dvipsnames]{xcolor}
\usepackage[margin=0.5in,hmarginratio=1:1]{geometry}% equal left and right margins
\usepackage{graphicx}
% \usepackage{fontspec}
% \setmainfont{Times New Roman}
\usepackage{titling}
% \usepackage{book-of-common-prayer}
\usepackage{yfonts}

\usepackage{afterpage}
\newcommand\myemptypage{
    \null
    \thispagestyle{empty}
    \addtocounter{page}{-1}
    \newpage
    }

\setlength{\droptitle}{-1in}

\def\ae{æ}
\def\AE{Æ}
\def\oe{œ}
\def\OE{Œ}


\begin{document}
\pagenumbering{roman}


\title{\includegraphics[width=2.5in]{Picture1.png} \\ \Huge \bfseries \color{Red} B R E V I A R V M \\ \huge R O M A N V M, \\ \large \mdseries \color{Black} Ex sacra potissimum scriptura, \& probatis sanctorum \\ historiis nuper consectum, ac denuo per \\ eundem authorem accurantius \\ recognitum. \\ \ \\ \itshape scrutamini scripturas, quoniam illæ sunt, \\ quæ testimonium perhibent de me. \\Iohan. V.}

\author{VENETIIS APVD IVNCTAS}
\date{MDCXXXI}

\maketitle
\bookmark[page=1]{TITULUM}
\myemptypage

% \frontmatter

\pagenumbering{roman}

\begin{center} \color{red} \hypertarget{Preface} \Large
AD SANCTISSIMVM PATREM \\ \& Dominum nostrum Paulum tertium \\ Pontificem Maximum,
\end{center}

\begin{center} \color{black} \large
Francisci Quignonij Tituli Sanctæ Crucis in Ierusalem presbyteri Cardinalis, in Breuiarium proxime consectum, ac denuo \\ recognitum.
\end{center}
\vspace{+.25em}
\centerline{\Large \color{red} P R \AE \ F A T I O.}
\bookmark[dest=Preface]{PR\AE FATIO}

\Huge \textbf{B}
\normalsize Reviarum Romanum nuper a nobis felicis recordationis Clementis. VII. Pontifi. Maximi hortatu consectum, ac potius in ampliorem sacrarum scripturam lectionem, ad veterem sanctorum patrum, \& conciliorum antiquorum formam reuocatum, tuaque voluntate sanctissime pater editum, graues plerosque ac doctos viros ita probasse \& recepisse intellexi, vt nihil in eo metandum existimarent.
Alios item animaduerti graues etiam, \& prudentes homines, qui eius rationem magnopere probantes, nonnihil tamen in eo desiderari affimarent.
Illud vero nunquam dubitatui, fore in tanta multitudine nonnullos ex ijs videlicet, qui in diuerso precandi ritu consenuissent, quibus labor ille noster non esset perinde gratus, existimantibus ab inueterata illa consuetudine precandi nulla ratione clericis esse discedendum.
Imo vero nobis primam editionem Breuiarij, non tanquam promulgationem legis esse placuerat, sed quasi publicam quandam deliberationem:
vt sic proposita nostra sententia, iudicia multorum exquireremus, \& quod omnium commodissimum, \& religioni, ac pietati conuenientissimum plerisque prudentibus, grauibusque viris visum esset, sequeremur.
Nihil enim humano elaboratum ingenio tam exactum initio vnquam fuit, quin postea multorum accedente iudicio perfectius reddi possit, vt in ipsis etiam ecclesiasticis institutis, circa primitiuam presertim ecclesiam contigisse videmus.
Itaque, multorum sententiis 
% end pp7
collatis, qu\ae \ nobis partim vocibus, partim scriptis innotuerunt, iudicium eorum secuti, qui omnium prudentissime sentire visi sunt, libenter quedam addidimus, alia mutauimus, \& omnia diligenter recongnouimus, retenta tamen summa forma Breuiarii.
Sed quoniam sic fert natura rerum, vt nihil sit tam rectum, nihil ranta ratione in vitam vsumque hominum inductum, cuius nouitas non sit aliquibus ingrata, non temere facturi esse videmur, si rationem totius instituti nostri a nobis prius summatim redditam, nunc accuratius recognito Breuiario, paulo latius explicabimus.
Mihi enim (vt s\ae pe sum professus) cogitanti, atque animo repetenti initia veteris insituti, quo sancitum est, vt clerici sacris initiati, vel sacerdotiis pr\ae sidentes, singulis diebus perlegant horarias preces, quas Canonicas etiam appellamus, tres pr\ae cipue caus\ae \ spectat\ae \ fuisse videri solent.
Quarum ea prima est, qui cum c\ae teri homines in quaque ciuitate, aut suum quisque negocium agant, aut in Republica administranda sint occupati, clericis ex eo vocatis, qui sortis Domini sint, quique bonis ecclesiasticis aluntur, hoc potissimum negocium diuinis, \& humanis legibus est iniunctum, vt Deum habere propitium in commissum sibi populum, de se bene merentem cunctis rationibus enitantur.
quod %sic
non solum sacrificiis efficitur, sed etiam precibus, qu\ae \ a pio corde proficiscantur.
ipsi %sic
enim (vt testimonio versibusque vtar beati Alexandri pap\ae \ \& martyris) pro populo interpellant, \& populi peccata comedunt, quia precibus suis \& oblationibus ea delent, atque consumunt.
Qui quanto digniores fuerint, tanto facilius pro quibus necessitatibus clamant, exaudiuntur.
Quod testimonio probat Iacobi Apostoli, qui nos ad precandum cohortans, Orate (inquit) pro inuicem vt saluemini.
multum %sic
enim valet deprecatio iusti assidua.
Altera causa est, vt qui reliquo populo exemplo debent esse virtutis \& sanctimoni\ae , assidua precatione Deum alloquentes, minus opportuni reddantur tentatori diabolo, si eos inuenerit, vt Hieronymus ait, occupatos, \& a cogitationibus caducarum rerum subinde auocati, contemplationi diuinarum assuescant.
Tertia, vt religionis quoque futuri magistri quotidiana sacr\ae \ scriptur\ae \ \& ecclesiasticum historiarum lectione erudiantur, complectanturque (vt Paulos ait) eum, qui secundum doctrinam est, fidelem 
%end pp8
sermonem, \& potentes sint exhortari in doctrina sana, \& eos, qui contradicunt, arguere.
Et profecto si quis modum precandi olim a maioribus traditum diligenter confideret, plane inelligat horum omnium pr\ae cipia ab ipsis habitam esse rationem.
Sed factum est nescio quo pacto precantium negligentia, vt paulatim a sanctissimis illis veterum patrum institutis discederetur.
Nam libri scripure sacr\ae, qui statis anni temporibus legendi erant more maiorum (vt est in monumentis Gelasii pap\ae, \& concilii Romani. lxx. episcoporum) vixdum incepti omittuntur in alio Breuiario.
Tam histori\ae \ sanctorum quedam tam inculte, \& tam sine delectu script\ae \ habentur in eodem, vt nec authoritatem habere videantur, nec grauitatem.
Accedit tam perplexus ordo, tamque difficilis precandi ratio, vt interdum paulo minor opera in requirendo ponatur, qui, cum inueneris, in legendo.
Quibus rebus animaduersis, felicis recordationis Clemens. VII. Pontifex Maxi. cum intelligeret officii sui esse, cum aliorum Christianorum commoditatibus prospicere, tum in primis clericorum, quibus ministris vteretur in commisso sibi grege administrando, me hortatus est, negociumque dedit, vt quantum cura, \& diligentia niti possem, preces horarias ita disponerem, \& ad veterm illam formam, quoad fieri posset, redigerem:
vt difficultatibus, dispendiisque sublatis, nec a summa rationis olim ab antiquis sanctisque patribus institut\ae, discederetur:
nec rursum clerici ab officio precandi deterrerentur laboris magnitudine.
Quam ego pro uinciam libentissime suscepi, simul vt Christi vicario honestissima precipienti parerem, simul vt bono publico pro mea virili parte seruirem.
Adhibitis igitur quibusdam meorum domesticorum prudentibus hominibus, ac sacrarum literarum \& pontificii iuris doctrina pr\ae ditis, eisdemque Gr\ae ce \& Latine eruditis, dedi operam qui maxime potui, vt commode ac ex vtilitate publica rem conficerem.
Ac illud ante omnia visum nobis est in consuetudinem reuocare, vt scriptura sacra omnia visum nobis est in consuetudinem reuocare, vt scriptura sacra maxime omnium toto anno, \& omnes %oes
psalmi singulis septimanis perlegerent.
Utrumque enim veteribus patribus summa ratione placuisse videbamus, qui plerosque libros (vt supra dictum est) utrumsque testamenti per anni tempora legendos disposuerant in Breuiario, \& per dies sic psalmos distribuerant, vt singulis septimanis
%end pp9
totum psalterium legeretur, quod singulis diebus in primitua ecclesia perlegi solitum esse tradit Speculator.
Sed (vt supra diximus) precantium negligentia factum erat, vt minima pars scriptur\ae \ sacr\ae \ legeretur, \& psalmorum plerisque omissis, pauci singulis fere diebus repeterentur.
Quod eis accidit, dum officij ferialis optime sanctissimeque, a maioribus instituti longitudinem \& laborem refugerent.
Ad quem vitandum, varia compendia excogitantes, consuetudinem induxerunt omissis feriis, toto fere anno scriptorum %sctoru
festa celebrandi, etiam in quadragesima.
quod %sic
quantum sit contra veterum instituta, declarat concilium Laodicenum, quo cautum est, ne in quadragesima martyrum natalitia celebrentur, \& item Toletanum decimum, in quo traditur, nihil esse in quadragesima de scriptorum %sctoru
solennitatibus %sic
celebrandum, idque de antiquitate regulari cautum esse.
Quod ipsi animaduertentes, conditionemque temporum \& infirmitatem clericorum spectantes, rationem ineundam putauimus, vt pari propemodum labore, vt cum festum celebrabatur, sed minor qui cum agebatur de feria veteris testamenti, magna \& vtilissima pars, \& totum nouum, pr\ae ter partem Apocalypsis, in anno, \& singulis Septimanis psalmi omnes perlegerentur, terni singulis horis, vnius longitudine cum alterius breuitate sic compensata, vt labor legendi diurnus par propemodum sit omnibus diebus, siue feria sit, siue festum celebretur, habita quoque in distribuendo, quoad fieri potuit, ratione feriarum, festorum non perinde, nec enim si cuius psalmi vnum aut alterum verbum potest ad aliquod festum accommodori, idcirco mutandus erat ordo commodissimus, quo nullus psalmorum in hebdomada pr\ae termittitur, vt non illa verba solum, sed c\ae tera omnia in ipsis contenta mysteria, liceat intelligentibus septimoquoque die contemplari.
Nam ad psalterium perlegendum, si veterem illam rectam \& sanctam psalmorum distributionem ferialem, in qua psalmi quidam quotidie repetuntur, retinuissemus, longitudo officij perinde vt antea clericos deterreret, qui hac commodissima, \& ad vetusinstitutum %not sure
afformata mediocritate, multum, vtiam experi sumus, alliciuntur.
Nam licet opera consueta legendi scripturas, parum minuatur in hoc Breuiario, presertim cum pleraque omnia ex libro legi sit necesse, curauimus tamen, vt ordine longe simpliciore, \& expeditiore magna 
%end page 10
pars laboris minueretur, sublato t\ae dio \& difficultate, psalmos, \& alia hincinde cum mora \& molestia perquirendi, \& difficillimam rationem longo tempore perdiscendi.
Versiculos, responsoria, \& capitula omitiere %omittere??
idcirco visum est, non qui h\ae c superuacanea aut inutilia viderentur (augent enim pietatem, \& sunt scriptur\ae \ sacr\ae \ particul\ae ) sed quam %qm
cum introduct\ae \ sint ad cantus potissimum modulandes, \& legentes s\ae pe morentur cum molestia qu\ae ritandi, locum relinqui voluimus continenti lectioni scriptur\ae \ sacr\ae , qu\ae \ magis facit ad pietatem atque doctrinam priuatim legentium, quorum hac quidem in parte maiorem rationem haberi voluimus, qui concinentium.
Nec enim ad precandum cuncta vtilia \& salubria congeri debent, ne clerici grauentur iniquiori pondere.
Atque vtinam tam robusti essemus, vt totum etiam vetus testamentum in anno perlegere nongrauaremur, libenter enim omnes %oes
illius libros legendos proponeremus:
sed (vt dictum est) habenda fuit ratio, ne precandi labor in tanta clericorum infirmitate modum excederet.
Officium beat\ae \ virginis in alio Breuiario vsitatum omittitur, nec propterea minuitur ipsius honor.
Nam paucis festis exceptis, quotidie variis commemorationibus adoratur, \& obsecratur:
atque eius memoria multis dicatis eidem festis, \& octauis pro cuiusque ratione celebratur, ipsisque solidum officium tribuitur sabbatis fere omnibus pr\ae ter qui in Quadragesima, in qua ne ipsum quidem Annuntiationis festum celebrari placuit veteribus patribus:
vt docet idem Concilium, quod supra diximus, Toletanum.
Et profecto quorundam Psalmorum grauem plerisque repetitionem omitti, non iam molestum esse virgini matri credi par est, qui gratum illud, quod clerici ad ipsius filij Iesu Christi diuinum cultum commodissima \& expedita ratione alliciantur.
Ex scriptorum %sctoru
historiis diligenter consideratis, quasdam delegimus, ex probatis authoribus Gr\ae cis \& Latinis fere decerptas, \& has cultiore stylo non tamen fucato, sed potius ex industria \& iussu Clementis pontificis attemperato, %sic
disposuimus per annum totum.
Nam licet veteribus in quadragesima nulla (vt docuimus) scriptorum %sctoru
festa celebrari placuerit, tamen a iunioribus tam assidue in ea quoque celebrantur, vt a nobis ne tum quidem omnino festa saltem solennia sine aliquo pusillorum offendiculo vitari
%end pp11
posse putauerimus.
Vt igitur quasdam sanctorum historias retinuimus, aut reposuimus denuo conquisitas, sic quasdam omisimus illis nec probabilitate, nec grauitate pares, temeritate videlicet priuata quorundam introductas, qui dum fugerent officii ferialis longitudinem, festa scriptorum %sctoru
studiose conquirentes, quasdam in Breuiario pr\ae ter publicam authoritatem inconsultis Pontificibus inserendas curauerunt, quas Pontifices ipsi haud dubie in vsum ecclesi\ae \ non admisissent.
Sunt enim nonnull\ae \ (nisi fallimur) ex earum genere, que publico sanctissimoque decreto ab ecclesia reiiciuntur, \& quidem in eo loco siue capite celeberrimo in quo sancta Romana ecclesia quas scripturas clericis legendas admittat, \& rursus quas reiiciat, explicatur.
Vbi cum multa summo iudicio sancita, tum illud habetur sanctissime \& grauissime a sanctis patribus institutum de historiis scriptorum, %sctoru
quarum authores ignorantur.
Ibi enim Gelasius papa ex synodo septuaginta episcoporum, Gesta (inquit) scriptorum %sctoru
martyrum, qui multiplicibus tormentorum cruciatibus, \& mirabilibus confessionum triumphis irradiant, quid ita esse catholicorum dubitet?
\& maiora eos in agone fuisse perpessos?
nec suis viribus, sed Dei gratia \& antiquam \& consuetudinem in sancta Romana ecclesia non leguntur:
quia eorum, qui scripsere, nomina penitus ignorantur, \& ab infidelibus, aut dictis superflua, aut minus apta qui rei ordo fuerit scripta esse putantur.
Hactenus Gelasius papa.
Qui statim subiiciens exempla quorundam scriptorum, %sctoru
quorum histori\ae \ ab h\ae reticis probabantur esse conscript\ae , rursus ait:
Propter quod, vt dictum est, ne vel leuius subsannandi oriretur occasio, in sancta Romana ecclesia non leguntur, sic Gelasius.
Contra h\ae c igitur, qu\ae \ summa ratione fuerant a sanctis patribus instituta, animaduertentes quedam, %quesdam?
temeritate quorundam priuata irrepsisse, diligenter curauimus, vt qu\ae \ publico ecclesi\ae \ decreto reiiciuntur, omnino vitaremus, \& in quibusdam aliis resecatis, qu\ae \ ad rem minus pertinerent, probatissima qu\ae que consectaremur, iuxta pr\ae ceptum Pauli, quem his quoque in rebus authorem adhibere iubemur ibidem ab ecclesia.
Gelasius enim in eodem loco de nonnullis historiis certitudinis non admodum explorat\ae , qu\ae \ a catholicis legerentur,
%end pp12
loquutus, h\ae c subiicit.
Sed cum hoc ad catholicorum manus peruenerit, beati Pauli Apostoli sententia pr\ae cedat:
Omnia probate:
quod bonum est tenete.
Precum pro Pontificibus \& regibus, \& fidelium bonis spiritalibus ac temporalibus, diligenter est adhibita ratio.
Nam singulis diebus ad Matutinum \& Vesperam, nostris temporibus pacem, ecclesi\ae \ iustitiam \& sanctitatem precamur, \& fideles omnes %oes
tum viuentes, tum etiam defunctos Deo commendamus.
Quod impensius etiam efficitur, cum officium defunctorum, sive psalmi p\oe nitentiales cum litania \& orationibus leguntur, vt s\ae pe fit in quadragesima, quod prius tunc fiebat duntaxat cum ageretur de feria.
Sed hoc tam rarum erat consuetudine clericorum festa pro feriis celebratium, vt plerique ipsorum vix semel aut iterum toto anno legerent officium feriale.
Itaque si quis diligenter animaduerterit, \& vetus patrum concilium institutumque considerauerit, plane intelliget, hoc Breuiarium non tam esse nouum inuentum, qui Breuiarij antiqui adhibito quodam temperamento restitutionem.
Habet igitur h\ae c precandi ratio tres maximas commoditates.
Primam, qui precantibus simul acquiritur vtriumsque testamenti peritia.
Secundam, qui res est expeditissima propter summam ordinis simplicitatem, \& nonnullam breuitatem.
Tertiam, qui histori\ae \ scriptorum %sctoru
sic conscript\ae \ sunt, vt nihil habeant, quod graues \& doctas aures offendat.
Si cui ergo laboriosum in hoc Breuiario videbitur, pleraque omnia ex libro legi, cum multa in alio, qu\ae \ propter frequentem repetitionem ediscuntur, memoriter pronuntientur, compenset cum hoc labore cognitionem sacr\ae \ scriptur\ae , qu\ae \ sic in dies augescit, \& intentionem animi, quam Deus ante omnia in precantibus requirit (hanc enim maiorem legentibus, qui memoriter proferentibus adesse necesse est) \& huiusmodi laborem non modo fructuosum, sed etiam salutarem iudicabit.
Orandum est enim, vt Paulus ait, spiritu, \& mente.
Nam sine fructu mens eius est, qui orat lingua tantum, id est qui sine attentione, ac velut aliud agens precatur, quod fere accidit iis, qui psalmos memoriter percurrunt. Habes, pater sanctissime, instituti nostri rationem, habes formulam Breuiarij.
Superest vt quemadmodum c\oe pisit, aspires c\oe ptis nostris, labori faueas, public\ae \ commoditatis causa suscepto. \indent Vale.
%end pp13

\newpage

% pp14
\begin{center} \color{red} \hypertarget{Letter} \Large
Tabula literarum Dominicalium. 1564.
\end{center}
\bookmark[dest=Letter]{TABULA LITERARUM DOMINICALIUM}

\begin{center}
\begin{tabular}{c c c c c c c c c c c c c c}
b & g & f & e & d & b & \color{red} A & g & f & d & c & b & \color{red} A & f\\
\color{red} A & \  & \  & \  & c & \  & \  & \  & e & \  & \  & \  & g & \ \\
e & d & c & \color{red} A & g & f & e & c & b & \color{red} A & g & e & d & c\\
\  & \  & b & \  & \ & \  & d & \  & \  & \  & f & \  & \  & \ \\
\end{tabular}
\end{center}

\Huge \textbf{\color{red} I}
\normalsize N Tabula suprascripta, literam Dominicalem singulis annis hoc pacto inuenius.
Primam literam, qu\ae \ est \color{red} b a \color{black} tribue anno. \color{red} 1564.
\color{black} Sequentem, qu\ae \ est \color{red} g \color{black} tribue anno. \color{red} 1565.
\color{black} Sicque deinceps progredere quousque peruenias ad numerm anni, cuius literam scire vis.
Quia litera super quam cadit numerus illius anni, erit litera Dominicalis.
Qu\ae \ si vnica occurrerit, annum communem esse intelliges:
si duplex, Bissextilem.
Et tunc prior, id est superior, usque ad festim Mathi\ae \ apostoli deseruiet.
Inferior autem ad reliquam anni partem accommodabitur.
Sciendum autem qui finitis \color{red} 18. \color{black} annis, redeundum est ad principium cum annis tunc currentibus.

\begin{center} \color{red} \hypertarget{Golden} \Large
Aureus numeris.
\end{center}
\vspace{-1em}
\begin{center}
7.8.9.10.11.12.13.14.15.16.17.18.19.1.2.3.4.5.6.7.8.\\
9.10.11.12.13.14.15.
\end{center}
\bookmark[rellevel=1,dest=Golden]{AUREUS NUMERIS}

\Huge \textbf{\color{red} S}
\normalsize Imili modo, vt dictum est de litera Dominicali deprehendes etiam aureum numerum in suprascripta numerorum serie.
Nam primum numerum, qui est. \color{red} 7. \color{black} tribue anno. \color{red} 1564.
\color{black} Sequentem qui est. \color{red} 8. \color{black} tribue anno. \color{red} 1565. \color{black} \& sic deinceps, donec ad annum tuum peruenias: \& finitis. \color{red} 19. \color{black} annis, regredere iterum ad principium.

\begin{center} \color{red} \hypertarget{NewMoon} \Large
Nouiluniorum Canon.
\end{center}
\bookmark[dest=NewMoon]{NOUILUNIORUM CANON}
\vspace{-1em}
\Huge \textbf{\color{red} N}
\normalsize Ouilunium singulis mensibus sic reperitur:
Habito aureo numero anni tui, vbicunque ille in ordine aurei numeri in Calen. posito tibi occurrerit, a loco illius diei iuxta positi sequentium verborum syllabas sursum computa.
In c\oe lis est \color{red} hic. \color{black} vnicuique diei accommodando syllabam vnam.
Diesque, ille supra quem syllaba, \color{red} hic. \color{black} ceciderit, dies est Nouilunii.
\color{red} I\color{black} n sequenti Calend. cuilibet festo habenti aliquid proprium apponitur, quoto folio id inueniatur.
Alia vero festa, quibus non apponitur folio % f'.
nihil habent proprium.
Item in quibusdam diebus assignantur terti\ae \ lectiones repetend\ae \ ex epistolis cum fit officium de feria, quarum nonnull\ae \ in quibusdam annis omittuntur, in aliis autem leguntur, vt infra in regulis generalibus videbitur.


\newpage

%pp15
\begin{center} \hypertarget{Kalendar} \Large
IANVARIVS HABET DIES XXXI.
\end{center}
\bookmark[startatroot,dest=Kalendar]{KALENDARIUM}
\bookmark[rellevel=1,dest=Kalendar]{IANVARIVS}

\begin{center}
\begin{tabular}{l r l r r l}
\multicolumn{2}{l}{\color{red}Au. nu.} & & & \color{red} Dies & \\
iij & \color{red} A & \multicolumn{2}{c}{\color{red} Calendis} & 1 & \color{red}Circuncisio Christi duplex maius. \color{black} \#\#.\\
 & b & iiij & no. & 2 & Octaua sancti Stephani duplex mi. \color{black} \#\#.\\
xj & c & iij & no. & 3 & Octaua sancti Ioannis euan. du. mi. \color{black} \#\#.\\
 & d & \color{red} prid. & no. & 4 & Octaua sanctorum Innocen. dup. mi. \color{black} \#\#.\\
xix & e & \multicolumn{2}{c}{\color{red} Nonis} & 5 & De vigilia epiphani\ae . \color{black} \#\#.\\
viij & f & viij & id. & 6 & \color{red} Epiphania Christi duplex maius. \color{black} \#\#.\\
 & g & vij & id. & 7 & De octaua Epiphani\ae . \color{black} \#\#.\\
xvj & \color{red} A & vj & id. & 8 & De octaua Epiphani\ae . \color{black} \#\#.\\
v & b & v & id. & 9 & De octaua Epiphani\ae . \color{black} \#\#.\\
 & c & iiij & id. & 10 & De octaua Epiphani\ae . \color{black} \#\#.\\
xiij & d & iij & id. & 11 & De octaua Epiphani\ae . \color{black} \#\#.\\
ij & e & \color{red} pridi. & id. & 12 & De octaua Epiphani\ae . \color{black} \#\#.\\
 & f & \multicolumn{2}{c}{\color{red} Idibus} & 13 & De octaua Epiphani\ae \ duplex mi. \color{black} \#\#.\\
x & g & xix & cal. & 14 & Basilius episcopus confes. \color{red} Fuit. 1. hu. \color{black} \#\#.\\
 & \color{red} A & xviij & cal. & 15 & Martina virg: marty. \color{red} Fuit. 2. hu. \color{black} \#\#.\\
xviij & b & xvij & cal. & 16 & Marcellus papa martyr. \color{black} \#\#.\\
vij & c & xvj & cal. & 17 & Antonius abbas duplex minus. \color{black} \#\#.\\
 & d & xv & cal. & 18 & Cathedra Romana Petri dup. ma. \color{black} \#\#.\\
xv & e & xiiij & cal. & 19 & Telespho. papa mart. \color{red} Fuit. 5. hu. \color{black} \#\#.\\
iiij & f & xiij & cal. & 20 & \color{red} Fabianus papa \& Sebast. mar. du. mi. \color{black} \#\#.\\
 & g & xij & cal. & 21 & Agnes virgo mar. duplex minus. \color{black} \#\#.\\
xij & \color{red} A & xj & cal. & 22 & \color{red} Vincentius \& Anastasius mar. \color{black} \#\#.\\
j & b & x & cal. & 23 & Alphonsus archiepiscopus confes. \color{black} \#\#.\\
 & c & ix & cal. & 24 & Timotheus episcopus marty. \color{black} \#\#.\\
ix & d & viij & cal. & 25 & \color{red} Conuersio Pauli apostoli dup. ma. \color{black} \#\#.\\
& e & vij & cal. & 26 & Polycarpus episcopus marty. \color{black} \#\#.\\
xvij & f & vj & cal. & 27 & Ioannes Chrysostomus episcopus confes. \color{black} \#\#.\\
vj & g & v & cal. & 28 & Lucianus presby. mar. \color{red} Fuit. 7. hu. \color{black} \#\#.\\
 & \color{red} A & iiij & cal. & 29 & Paulus primus eremita. \color{red} Fuit. 10. hu. \color{black} \#\#.\\
xiiij & b & iij & cal. & 30 & Iginius papa mar. \color{red} Fuit. 11. hu. \color{black} \#\#.\\
 & c & \color{red} prid. & cal. & 31 & Hilarius episcopus confes. \color{red} Fuit. 13. hu. \color{black} \#\#.\\
\end{tabular}
\end{center}

%pp16
\newpage

\begin{center} \hypertarget{FEBRVA} \Large
FEBRVA. HABET DIES XXVIII.
\end{center}
\bookmark[dest=FEBRVA]{FEBRVARIUS}
\begin{center}
\color{red} Et si fuerit Bissextus. XXIX.
\end{center}

\begin{center}
\begin{tabular}{l r l r r l}
\multicolumn{2}{l}{\color{red}Au. nu.} & & & \color{red} Dies & \\
 & d & \multicolumn{2}{c}{\color{red} Calendis} & 1 & Ignatius episcopus martyr. folio. \color{black} \#\#.\\
xi & e & iiii & no. & 2 & \color{red} Purificatio Mari\ae \ virginis du. ma. \color{black} \#\#.\\
xix & f & iii & no. & 3 & Blasius episcopus martyr. \color{black} \#\#.\\
viii & g & \color{red} prid. & \color{red} no. & 4 & Phileas episcopus \& Philoromus mar. \color{black} \#\#.\\
 & \color{red} A & \multicolumn{2}{c}{\color{red} Nonis} & 5 & Agata virgo martyr. \color{black} \#\#.\\
xvi & b & viii & id. & 6 & Dorothea virgo martyr. \color{black} \#\#.\\
v & c & vii & id. & 7 & Adauctus, \& socii martyres. \color{black} \#\#.\\
 & d & vi & id. & 8 & Cointha virgo martyr. \color{black} \#\#.\\
xiii & e & v & id. & 9 & Apollonia virgo martyr. \color{black} \#\#.\\
ii & f & iiii & id. & 10 & Scholastica virgo. \\%sic
 & g & iii & id. & 11 & Prisca virgo mar. qu\ae \ fuit. 18 Ian. \color{black} \#\#.\\
x & \color{red} A & \color{red} prid. & \color{red} id. & 12 & Eulalia virgo martyr. \color{black} \#\#.\\
 & b & \multicolumn{2}{c}{\color{red} Idibus} & 13 & Ex epist. i. Petri. \color{red} P\color{black} ropter quod. \#\#.\\
xviii & c & xvi & cal. & 14 & Valentinus presbyter martyr. \color{black} \#\#.\\
vii & d & xv & cal. & 15 & Faustinus, \& Iovita marty. \color{black} \#\#.\\
 & e & xiiii & cal. & 16 & Iuliana virgo martyr. \color{black} \#\#.\\
xv & f & xiii & cal. & 17 & Ex epist. i. Petri. \color{red} D\color{black} eponentes. \#\#.\\
iiii & g & xii & cal. & 18 & Ex epist. i. Petri. \color{red} C\color{black} harissimi. \#\#.\\
 & \color{red} A & xi & cal. & 19 & Gabinus presbyter martyr. \color{black} \#\#.\\
xii & b & x & cal. & 20 & Ex epist. i. Petri. \color{red} S\color{black} imiliter. \#\#.\\
i & c & ix & cal. & 21 & Ex epist. i. Petri. \color{red} C\color{black} hristo igitur. \#\#.\\
 & d & viii & cal. & 22 & \color{red} Cathedra Antioch. Pet. apo. du. ma. \color{black} \#\#.\\
ix & e & vii & cal. & 23 & Ex epist. i. Petri. \color{red} S\color{black} eniores. \color{red} Vigilia. \color{black} \#\#.\\
 & f & vi & cal. & 24 & \color{red} Mathias aposto duplex maius. \color{black} \#\#.\\
xvii & g & v & cal. & 25 & Epist. Pauli ad Phil. \color{red} P\color{black} aulus. \#\#.\\
vi & \color{red} A & iiii & cal. & 26 & Ex epist. Pauli ad Philip. \color{red} M\color{black} ihi. \#\#.\\
 & b & iii & cal. & 27 & Iulianus, \& Eunus martyres. \color{black} \#\#.\\
xiiii & c & \color{red} prid. & \color{red} cal. & 28 & Ex epist. Pauli ad Philip. \color{red} I\color{black} taque. \#\#.\\
\end{tabular}
\end{center}


%pp17
\newpage

\begin{center} \hypertarget{MARTIVS} \Large
MARTIVS HABET DIES XXXI.
\end{center}
\vspace{+.25em}
\bookmark[dest=MARTIVS]{MARTIVS}

\begin{center}
\begin{tabular}{l r l r r l}
\multicolumn{2}{l}{\color{red}Au. nu.} & & & \color{red} Dies & \\
iii & d & \multicolumn{2}{c}{\color{red} Calendis} & 1 & Ex epist. Pauli ad Phil. \color{red} D\color{black} e c\ae te. fo. \#\#.\\%Ch3
 & e & iiii & no. & 2 & Ex epist. Pauli ad Philip. \color{red} I\color{black} taque fres. \#\#.\\%Ch4
xi & f & iii & no. & 3 & Emete. \& Cel. \& Aste. martyr. \color{black} \#\#.\\
 & g & \color{red} prid. & \color{red} no. & 4 & Lucius papa martyr. \color{black} \#\#.\\
xix & \color{red} A & \multicolumn{2}{c}{\color{red} Nonis} & 5 & Epistola ad Coloss. \color{red} P\color{black} aulus. \#\#.\\
viii & b & viii & id. & 6 & Epistola ad Coloss. \color{red} E\color{black} t vos. \#\#.\\
 & c & vii & id. & 7 & Thomas Aquin. confessor. \color{black} \#\#.\\
xvi & d & vi & id. & 8 & Epistola ad Coloss. \color{red} V\color{black} idete. \#\#.\\
v & e & v & id. & 9 & Quadraginta milites marty. \color{black} \#\#.\\
 & f & iiii & id. & 10 & \color{red} \AE quinoctium. \\%sic
xiii & g & iii & id. & 11 & \\%sic
ii & \color{red} A & \color{red} pridi. & \color{red} id. & 12 & Gregorius papa confes. duplex mi. \color{black} \#\#.\\
 & b & \multicolumn{2}{c}{\color{red} Idibus} & 13 & \\
x & c & xix & cal. & 14 & \\
 & d & xviii & cal. & 15 & \\
xviii & e & xvii & cal. & 16 & \\
vii & f & xvi & cal. & 17 & \\
 & g & xv & cal. & 18 & \\
xv & \color{red} A & xiiii & cal. & 19 & Ioseph confessor duplex minus. \#\#.\\
iiii & b & xiii & cal. & 20 & \\
 & c & xii & cal. & 21 & Benedictus abbas duplex minus. \#\#.\\
xii & d & xi & cal. & 22 & \\
i & e & x & cal. & 23 & \\
 & f & ix & cal. & 24 & \\
ix & g & viii & cal. & 25 & \color{red} Annuntiatio Mari\ae \ virgi. du. ma. \color{black} \#\#.\\
 & \color{red} A & vii & cal. & 26 & \\
xvii & b & vi & cal. & 27 & \\
vi & c & v & cal. & 28 & \\
 & d & iiii & cal. & 29 & \\
xiiii & e & iii & cal. & 30 & Ex epist. ad Coloss. \color{red} I\color{black} gitur si. \#\#.\\
iii & f & \color{red} prid. & \color{red} cal. & 31 & Ex epist. ad Coloss. \color{red} D\color{black} omini quod. \#\#.\\
\end{tabular}
\end{center}


%pp18
\newpage

\begin{center} \hypertarget{APRILIS} \Large
APRILIS HABET DIES XXX.
\end{center}
\vspace{+.25em}
\bookmark[dest=APRILIS]{APRILIS}

\begin{center}
\begin{tabular}{l r l r r l}
\multicolumn{2}{l}{\color{red}Au. nu.} & & & \color{red} Dies & \\
 & g & \multicolumn{2}{c}{\color{red} Calendis} & 1 & Ex epist. j. ad Thess. \color{red} P\color{black} aulus. folio. \#\#.\\%Ch1
xj & \color{red} A & iiij & no. & 2 & Ex epist. ad Thess. \color{red} I\color{black} deo \& nos. \#\#.\\%Ch2:13
 & b & iij & no. & 3 & Pancratius martyr. \color{black} \\
xix & c & \color{red} prid. & no. & 4 & Isidorus episcopus confessor. \color{black} \#\#.\\
viij & d & \multicolumn{2}{c}{\color{red} Nonis} & 5 & Vincentius confessor ordinis pr\ae di. \color{black} \\
xvj & e & viij & id. & 6 & Xystus papa martyr. \color{black} \#\#.\\
v & f & vij & id. & 7 & Ex epist. prima ad Thess. \color{red} D\color{black} e cetero. \#\#.\\
 & g & vj & id. & 8 & Dionysius episcopus confessor. \color{black} \\
xiij & \color{red} A & v & id. & 9 & Ex epist. prima ad Thes. \color{red} D\color{black} e tempo. \#\#.\\
ij & b & iiij & id. & 10 & Ex epist secunda ad Thess. \color{red} P\color{black} aulus. \#\#.\\
 & c & iij & id. & 11 & Leo primus papa confess. \color{black} \#\#.\\
x & d & \color{red} pridi. & id. & 12 & Ex epist. secunda ad Thes. \color{red} R\color{black} ogamus. \#\#.\\
 & e & \multicolumn{2}{c}{\color{red} Idibus} & 13 & Iustinus philosophus martyr \#\#.\\
xviij & f & xviij & cal. & 14 & Valerian. Tybur. \& c. martyres. \#\#.\\
vij & g & xvij & cal. & 15 & Ex epist. secunda ad Thes. \color{red} D\color{black} e cetero. \#\#.\\
 & \color{red} A & xvj & cal. & 16 & Ex epist ad Ephesi. \color{red} P\color{black} aulus. \#\#.\\
xv & b & xv & cal. & 17 & Anicetus papa martyr. \\
iiij & c & xiiij & cal. & 18 & Apollonius Senator marty. \#\#.\\
 & d & xiij & cal. & 19 & Ex epist ad Ephesi. \color{red} E\color{black} t vos. \#\#.\\
xij & e & xij & cal. & 20 & Ex epist. ad Ephes. \color{red} H\color{black} uius rei. \\
j & f & xj & cal. & 21 & Ex epist. ad Ephes. \color{red} O\color{black} bsecro. \#\#.\\
 & g & x & cal. & 22 & Caius papa martyr. \#\#.\\
ix & \color{red} A & ix & cal. & 23 & \color{red} Georgius martyr.\\
 & b & viij & cal. & 24 & Ex epist. ad Ephes. \color{red} R\color{black} enouamini. \#\#.\\
xvij & c & vij & cal. & 25 & \color{red} Marcus euangelista duplex maius. \color{black} \#\#.\\
vj & d & vj & cal. & 26 & Cletus \& Marcel. pap\ae , \& mar. \#\#.\\
 & e & v & cal. & 27 & Anastatius papa confessor.\\
xiiij & f & iiij & cal. & 28 & Vitalis martyr.\\
iij & g & iij & cal. & 29 & Petrus ordinis pr\ae dica. mart. \#\#.\\
 & \color{red} A & \color{red} prid. & cal. & 30 & Ex epistola ad Ephes. \color{red} V\color{black} idete. \#\#.\\

\end{tabular}
\end{center}



%pp19
\newpage

\begin{center} \hypertarget{MAIUS} \Large
MAIUS HABET DIES XXXI.
\end{center}
% \vspace{+.25em}
\bookmark[dest=MAIUS]{MAIUS}

\begin{center}
\begin{tabular}{l r l r r l}
\multicolumn{2}{l}{\color{red}Au. nu.} & & & \color{red} Dies & \\
xi & b & \multicolumn{2}{c}{\color{red} Calendis} & 1 & \color{red} Philippi \& Iacobi apost. du. ma. \color{black} folio. \#\#.\\
 & c & vi & no. & 2 & Athanaius episcopus confes. \color{black} \#\#.\\
xix & d & v & no. & 3 & \color{red} Inuentio sancte Crucis duplex ma. \color{black} \#\#.\\
viii & e & iiii & no. & 4 & Monica mater sancti. Augustini. \color{black} \\
 & f & iii & no. & 5 & Alexander papa mar. \color{red} Fuit. 3. huius. \color{black} \#\#.\\
xvi & g & \color{red} prid. & \color{red} no. & 6 & \color{red} Ioannis ante portam Latinam dup. mi. \color{black} \#\#.\\
v & \color{red} A & \multicolumn{2}{c}{\color{red} Nonis} & 7 & Ex epist. ad Ephes. \color{red} F\color{black} ilii obed. \#\#.\\
 & b & viii & id. & 8 & \color{red} Apparitio sanctus Michaelis sup. mi. \color{black} \#\#.\\%scu'
xiii & c & vii & id. & 9 & Grego. Nazianze. episcopus confes. \color{black} \#\#.\\
ii & d & vi & id. & 10 & Gordianus \& Epimachus marty. \color{black}\\
 & e & v & id. & 11 & Epist. Iacobi. \color{red} I\color{black} acobus. \#\#.\\
x & f & iiii & id. & 12 & Nereus, Achileus, \& Pancra. marty. \color{black}\\
 & g & iii & id. & 13 & Ex epist. Iacobi. \color{red} E\color{black} stote. \#\#.\\
xviii & \color{red} A & \color{red} pridi. & id. & 14 & Victor \& Corona martyres. \#\#.\\
vii & b & \multicolumn{2}{c}{\color{red} Idibus} & 15 & Ex epist. Iacobi. \color{red} Q\color{black} uicunque. \#\#.\\
 & c & xvii & cal. & 16 & Ex epist. Iacobi. \color{red} N\color{black} olite. \#\#.\\
xv & d & xvi & cal. & 17 & Ex epist. Iacobi. \color{red} V\color{black} nde bella. \#\#.\\
iiii & e & xv & cal. & 18 & Ex epist. Iacobi. \color{red} A\color{black} gite nunc. \#\#.\\
 & f & xiiii & cal. & 19 & Potentiana virgo. \color{black} \#\#.\\
xii & g & xiii & cal. & 20 & Bernardinus confessor. \color{black} \#\#.\\
i & \color{red} A & xii & cal. & 21 & Iuo presbyt. confes. \color{red} Fuit. 19. huius. \color{black} \\
 & b & xi & cal. & 22 & Ex epist. Philip. \color{red} P\color{black} aulus \& Timot. \#\#.\\
ix & c & x & cal. & 23 & Ex epist. Philip. \color{red} M\color{black} ihi enim. \#\#.\\
 & d & ix & cal. & 24 & Ex epist. Philip. \color{red} I\color{black} taque. \#\#.\\
xvii & e & viii & cal. & 25 & Vrbanus papa martyr. \color{black} \#\#.\\
vi & f & vii & cal. & 26 & Eleuterius papa martyr. \\
 & g & vi & cal. & 27 & Ioannes papa marty. \\
xiiii & \color{red} A & v & cal. & 28 & Germanus episcopus confessor. \\
iii & b & iiii & cal. & 29 & Ex epist. ad Philip. \color{red} D\color{black} e c\ae tero. \#\#.\\
 & c & iii & cal. & 30 & Felix papa marty. \color{black} \#\#.\\
xi & d & \color{red} prid. & cal. & 31 & Petronilla virgo. \color{black} \\
\end{tabular}
\end{center}



%pp20
\newpage

\begin{center} \hypertarget{IVNIVS} \Large
IVNIVS HABET DIES XXX.
\end{center}
% \vspace{+.25em}
\bookmark[dest=IVNIVS]{IVNIVS}

\begin{center}
\begin{tabular}{l r l r r l}
\multicolumn{2}{l}{\color{red}Au. nu.} & & & \color{red} Dies & \\
 & e & \multicolumn{2}{c}{\color{red} Calendis} & 1 & Pamphilus presbyter martyr. \color{black} folio. \#\#.\\
xix & f & iiii & no. & 2 & Marcellinus, \& Petrus martyres. \color{black} \#\#.\\
viii & g & iii & no. & 3 & Ex epist. ad Philip. \color{red} I\color{black} taque. \#\#.\\
xvi & \color{red} A & \color{red} prid. & \color{red} no. & 4 & Ex epist. i. Ioan. \color{red} Q\color{black} uod fuit. \#\#.\\
v & b & \multicolumn{2}{c}{\color{red} Nonis} & 5 & Ex epist. i. Ioan. \color{red} F\color{black} ilioli mei. \#\#.\\
 & c & viii & id. & 6 & Ex epist. i. Ioan. \color{red} F\color{black} ilioli nouiss. \#\#.\\
xiii & d & vii & id. & 7 & Ex epist. i. Ioan. \color{red} V\color{black} idete. \#\#.\\
ii & e & vi & id. & 8 & Ex epist. i. Ioan. \color{red} C\color{black} harissimi. \#\#.\\
 & f & v & id. & 9 & Primus, \& Felicianus martyres. \\
x & g & iiii & id. & 10 & Ex epist. i. Ioan. \color{red} O\color{black} mnisqui. \#\#.\\
 & \color{red} A & iii & id. & 11 & \color{red} Barnabas apostolus du. ma. Solstitium. \color{black} \#\#.\\
xviii & b & \color{red} pridi. & \color{red} id. & 12 & Basilides, Cyrinus, \& c. martyres. \color{black}\\
vii & c & \multicolumn{2}{c}{\color{red} Idibus} & 13 & Antonius confes. ordinis mino. \color{black} \#\#.\\
 & d & xviii & cal. & 14 & Ex epist. ii. Ioan. \color{red} S\color{black} enior elect. \#\#.\\
xv & e & xvii & cal. & 15 & Vitus, Modestus \& Crescen. martyres.\\
iiii & f & xvi & cal. & 16 & Ex epist. iii. Ioan. \color{red} S\color{black} enior Caio. \#\#.\\
 & g & xv & cal. & 17 & Ex epist. i. ad Coloss. \color{red} P\color{black} aulus. \#\#.\\
xii & \color{red} A & xiiii & cal. & 18 & Marcus, \& Marcellianus mar. \#\#.\\
i & b & xiii & cal. & 19 & Geruasius, \& Prothasius marty. \color{black} \#\#.\\
 & c & xii & cal. & 20 & Siluerius papa martyr. \\
ix & d & xi & cal. & 21 & Ex epist ad Coloss. \color{red} E\color{black} t vos cum. \#\#.\\
 & e & x & cal. & 22 & Paulinus episcopus confessor. \\
xvii & f & ix & cal. & 23 & Ex epist. ad Coloss. \color{red} V\color{black} idete. \color{red} Vigilia. \color{black} \#\#.\\
vi & g & viii & cal. & 24 & \color{red} Natiuitas sancti Ioan. Bap. du. ma. \color{black} \#\#.\\
 & \color{red} A & vii & cal. & 25 & De octaua sancti Ioannis. \color{black} \#\#.\\
xiiii & b & vi & cal. & 26 & \color{red} Ioannis \& Pauli mar. dup. mi. \color{black} \#\#.\\
iii & c & v & cal. & 27 & De octaua sancti Ioannis. \#\#.\\
 & d & iiii & cal. & 28 & De octaua sancti Ioannis. \color{red} Vigilia. \color{black} \#\#.\\
xi & e & iii & cal. & 29 & \color{red} Petri \& Pauli apostoli dup. ma. \color{black} \#\#.\\
 & f & \color{red} prid. & \color{red} cal. & 30 & \color{red} Commemo. Pauli apostoli dup. mi. \color{black} \#\#.\\
\end{tabular}
\end{center}



%pp21
\newpage

\begin{center} \hypertarget{IVLIVS} \Large
IVLIVS HABET DIES XXXI.
\end{center}
% \vspace{+.25em}
\bookmark[dest=IVLIVS]{IVLIVS}

\begin{center}
\begin{tabular}{l r l r r l}
\multicolumn{2}{l}{\color{red}Au. nu.} & & & \color{red} Dies & \\
xix & g & \multicolumn{2}{c}{\color{red} Calendis} & 1 & Octa. nati. sancti Ioan. Bap. du. mi. \color{black} folio. \#\#.\\
viii & \color{red} A & vi & no. & 2 & \color{red} Visitatio Mari\ae \ virginis dup. ma. \color{black} \#\#.\\
 & b & v & no. & 3 & De octaua Visitationis. \color{black} \#\#.\\
xvi & c & iiii & no. & 4 & De octaua Visitationis. \color{black} \#\#.\\
v & d & iii & no. & 5 & De octaua Visitationis \color{black} \#\#.\\
 & e & \color{red} prid. & no. & 6 & Osta. apost. Petri \& Pauli dup. mi. \color{black} \#\#.\\
xiii & f & \multicolumn{2}{c}{\color{red} Nonis} & 7 & De octaua Visitationis. \#\#.\\
ii & g & viii & id. & 8 & De octaua Visitationis \color{black} \#\#.\\
 & \color{red} A & vii & id. & 9 & Octaua Visitationis duplex mi. \color{black} \#\#.\\
x & b & vi & id. & 10 & Septem frates martyres. \color{black} \#\#.\\
 & c & v & id. & 11 & Pius papa confessor. \#\#.\\
xviii & d & iiii & id. & 12 & Nabor \& Felix martyres. \color{black}\\
vii & e & iii & id. & 13 & Anacletus papa martyr. \#\#.\\
 & f & \color{red} pridi. & id. & 14 & Processus \& Martini. \color{red} Fuit. 2. huius. \color{black} \#\#.\\
xv & g & \multicolumn{2}{c}{\color{red} Idibus} & 15 & Bonauen. card. episc. conf. \color{red} Fuit. 13. hu. \color{black} \#\#.\\
iiii & \color{red} A & xvii & cal. & 16 & Eustachius episcopus confessor. \\
 & b & xvi & cal. & 17 & Alexius confessor. \\
xii & c & xv & cal. & 18 & Symphorosa cum septem fil. mar. \#\#.\\
j & d & xiiii & cal. & 19 & Iusta \& Ruffina virgi. martyr. \\
 & e & xiii & cal. & 20 & Margareta virgo martyr. \\
ix & f & xii & cal. & 21 & Praxedis virgo. \#\#.\\
 & g & xi & cal. & 22 & \color{red} Maria Magdalena duplex mi. \color{black} \#\#.\\
xvii & \color{red} A & x & cal. & 23 & Apollinaris episcopus martyr. \\
vi & b & ix & cal. & 24 & Christina virgo martyr. \color{red} Vigilia. \color{black} \\
 & c & viii & cal. & 25 & \color{red} Iacobus apostolus duplex maius. \color{black} \#\#.\\
xiiii & d & vii & cal. & 26 & \color{red} Anna mater Mari\ae \ virgi. dup. mi. \color{black} \#\#.\\
iii & e & vi & cal. & 27 & Pantaleon martyr. \\
 & f & v & cal. & 28 & Nazarius \& socij martyr. \\
xi & g & iiii & cal. & 29 & Martha virgo. \#\#.\\
xix & \color{red} A & iii & cal. & 30 & Abdon \& Sennen martyres. \\
 & b & \color{red} prid. & cal. & 31 & Nemesius \& Lucilla ma. \color{red} Fuit. 35. hu. \color{black} \#\#.\\
\end{tabular}
\end{center}



%pp22
\newpage

\begin{center} \hypertarget{AVGVSTVS} \Large
AVGVSTVS HABET DIES XXXI.
\end{center}
% \vspace{+.25em}
\bookmark[dest=AVGVSTVS]{AVGVSTVS}

\begin{center}
\begin{tabular}{l r l r r l}
\multicolumn{2}{l}{\color{red}Au. nu.} & & & \color{red} Dies & \\
 & c & \multicolumn{2}{c}{\color{red} Calendis} & 1 & \color{red} Vincula sancti Petri duplex mi. \color{black} \#\#.\\
xvi & d & iiii & no. & 2 & Stephanus papa martyr. \color{black} \#\#.\\
v & e & iii & no. & 3 & \color{red} Inuentio sancti Stepha. protomar. \color{black} \#\#.\\
 & f & \color{red} prid. & no. & 4 & Dominicus confes. deplex mi. \color{black} \#\#.\\
xiii & g & \multicolumn{2}{c}{\color{red} Nonis} & 5 & Sanct\ae \ Mari\ae \ ad niues duplex mi. \color{black} \#\#.\\
ii & \color{red} A & viii & id. & 6 & \color{red} Transfiguratio Domini dup. ma. \color{black} \#\#.\\
 & b & vii & id. & 7 & Iustinus presbyter mar. \color{red} Fuit. 4. huius. \color{black} \#\#.\\
x & c & vi & id. & 8 & Cyriacus, Largus, \& c. martyres. \\
 & d & v & id. & 9 & Xystus. ij. pa. ma. \color{red} Fuit. 6. hu. Vigil. \color{black} \#\#.\\
xviii & e & iiii & id. & 10 & \color{red} Laurentius martyr duplex ma. \color{black} \#\#.\\
vii & f & iii & id. & 11 & De octaua sancti Laurentij. \color{black} \#\#.\\
 & g & \color{red} pridi. & id. & 12 & Clara virgo duplex minus. \color{black} \#\#. \\
xv & \color{red} A & \multicolumn{2}{c}{\color{red} Idibus} & 13 & De octaua sancti Laurentij. \#\#.\\
iiii & b & xix & cal. & 14 & De octaua sancti Laurentij. \color{red} Vigilia. \#\#.\\%scti
 & c & xviii & cal. & 15 & \color{red} Assumptio Mari\ae \ virginis du. ma. \color{black} \#\#.\\
xii & d & xvii & cal. & 16 & De octaua Assumptionis. \#\#.\\
i & e & xvi & cal. & 17 & Octaua sancti Laurentij dup. mi. \#\#.\\
 & f & xv & cal. & 18 & De octaua Assumptionis. \#\#.\\
ix & g & xiiii & cal. & 19 & De octaua Assumptionis. \color{black} \#\#.\\

xii & g & xiii & cal. & 20 & Bernardinus confessor. \color{black} \#\#.\\

i & \color{red} A & xii & cal. & 21 & Iuo presbyt. confes. \color{red} Fuit. 19. huius. \color{black} \\

 & b & xi & cal. & 22 & Ex epist. Philip. \color{red} P\color{black} aulus \& Timot. \#\#.\\

ix & c & x & cal. & 23 & Ex epist. Philip. \color{red} M\color{black} ihi enim. \#\#.\\

 & d & ix & cal. & 24 & Ex epist. Philip. \color{red} I\color{black} taque. \#\#.\\

xvii & e & viii & cal. & 25 & Vrbanus papa martyr. \color{black} \#\#.\\

vi & f & vii & cal. & 26 & Eleuterius papa martyr. \\

 & g & vi & cal. & 27 & Ioannes papa marty. \\

xiiii & \color{red} A & v & cal. & 28 & Germanus episcopus confessor. \\

iii & b & iiii & cal. & 29 & Ex epist. ad Philip. \color{red} D\color{black} e c\ae tero. \#\#.\\

 & c & iii & cal. & 30 & Felix papa marty. \color{black} \#\#.\\

xi & d & \color{red} prid. & cal. & 31 & Petronilla virgo. \color{black} \\

\end{tabular}

\end{center}



%pp23
\newpage

\begin{center} \hypertarget{SEPTEMBER} \Large
SEPTEMBER HABET DIES XXX.
\end{center}
% \vspace{+.25em}
\bookmark[dest=SEPTEMBER]{SEPTEMBER}

\begin{center}
\begin{tabular}{l r l r r l}
\multicolumn{2}{l}{\color{red}Au. nu.} & & & \color{red} Dies & \\

xi & b & \multicolumn{2}{c}{\color{red} Calendis} & 1 & nooooooooooooooooooooo \color{black} folio. \#\#
.\\

 & c & vi & no. & 2 & Athanaius episcopus confes. \color{black} \#\#.\\

xix & d & v & no. & 3 & \color{red} Inuentio sancte Crucis duplex ma. \color{black} \#\#.\\

viii & e & iiii & no. & 4 & Monica mater sancti. Augustini. \color{black} \\

 & f & iii & no. & 5 & Alexander papa mar. \color{red} Fuit. 3. huius. \color{black} \#\#.\\

xvi & g & \color{red} prid. & \color{red} no. & 6 & \color{red} Ioannis ante portam Latinam dup. mi. \color{black} \#\#.\\

v & \color{red} A & \multicolumn{2}{c}{\color{red} Nonis} & 7 & Ex epist. ad Ephes. \color{red} F\color{black} ilii obed. \#\#.
\\

 & b & viii & id. & 8 & \color{red} Apparitio sanctus Michaelis sup. mi. \color{black} \#\#.\\%scu'

xiii & c & vii & id. & 9 & Grego. Nazianze. episcopus confes. \color{black} \#\#.\\

ii & f & vi & id. & 10 & Gordianus \& Epimachus marty. \color{black}\\

 & e & v & id. & 11 & Epist. Iacobi. \color{red} I\color{black} acobus. \#\#.\\

x & f & iiii & id. & 12 & Nereus, Achileus, \& Pancra. marty. \color{black}\\

 & g & iii & id. & 13 & Ex epist. Iacobi. \color{red} E\color{black} stote. \#\#.\\

xviii & \color{red} A & \color{red} pridi. & id. & 14 & Victor \& Corona martyres. \#\#.\\

vii & b & \multicolumn{2}{c}{\color{red} Idibus} & 15 & Ex epist. Iacobi. \color{red} Q\color{black} uicunque. \#\#.\\

 & c & xvii & cal. & 16 & Ex epist. Iacobi. \color{red} N\color{black} olite. \#\#.\\

xv & d & xvi & cal. & 17 & Ex epist. Iacobi. \color{red} V\color{black} nde bella. \#\#.\\

iiii & e & xv & cal. & 18 & Ex epist. Iacobi. \color{red} A\color{black} gite nunc. \#\#.\\

 & f & xiiii & cal. & 19 & Potentiana virgo. \color{black} \#\#.\\

xii & g & xiii & cal. & 20 & Bernardinus confessor. \color{black} \#\#.\\

i & \color{red} A & xii & cal. & 21 & Iuo presbyt. confes. \color{red} Fuit. 19. huius. \color{black} \\

 & b & xi & cal. & 22 & Ex epist. Philip. \color{red} P\color{black} aulus \& Timot. \#\#.\\

ix & c & x & cal. & 23 & Ex epist. Philip. \color{red} M\color{black} ihi enim. \#\#.\\

 & d & ix & cal. & 24 & Ex epist. Philip. \color{red} I\color{black} taque. \#\#.\\

xvii & e & viii & cal. & 25 & Vrbanus papa martyr. \color{black} \#\#.\\

vi & f & vii & cal. & 26 & Eleuterius papa martyr. \\

 & g & vi & cal. & 27 & Ioannes papa marty. \\

xiiii & \color{red} A & v & cal. & 28 & Germanus episcopus confessor. \\

iii & b & iiii & cal. & 29 & Ex epist. ad Philip. \color{red} D\color{black} e c\ae tero. \#\#.\\

 & c & iii & cal. & 30 & Felix papa marty. \color{black} \#\#.\\

xi & d & \color{red} prid. & cal. & 31 & Petronilla virgo. \color{black} \\

\end{tabular}

\end{center}

















%pp27
\newpage

\begin{center} \color{red} \hypertarget{Canon} \Large
Canon sequentis Tabul\ae \ ad festa mobilia\\ perpetuo inuenienda.
\end{center}
\bookmark[startatroot,dest=Canon]{CANON SEQUENTIS TABUL\AE}

\Huge \textbf{I}
\normalsize N tabula infrascripta sic festa mobilia singulis annis inuenies. Primo quere aureum numerum anni cuius efsta perquiris, deinde sub eius linea quere etiam literam Dominicalem eiusdem anni, qu\ae \ primo occurrerit:
\& notanter dicitur sub linea:
nam litera Dominicalis qu\ae \ e regione aurei numeri fuerit, eidem numero non inseruit, sed solum superioribus.
Postea in eadem linea transuersali, in qua literam Dominicalem inuenisti, inuenies etiam quo mense, \& quo die singula festa mobilia incidant, \& quot Dominic\ae \ Vagantes sint ante Septuagesimam, omnia h\ae c sub suis titulis.
In anno autem bissextili quot Dominic\ae \ Vagantes occurrant ante Septuagesimam, item quo die Septuagesima ipsa incidat, \& quo Feria quarta Cinerum, inuenies per priorem eius anni literam, qu\ae \ inseruit vsque ad festum Mathi\ae .
Reliqua vero omnia per posteriorem exquiruntur.
Sed has duas literas anni bissextilis sic accipere oportet, vt neutra earum fit e regione aurei numeri, nam tunc omissis illis, recurrendum est ad alias similes inferius %inserius?
primo loco occurrentes.
Pr\ae terea si in anno bissextili dies Cineram in mense Martio incidisse supradicta ratione deprehendatur, per vnum diem retrahenda erit.
In proximis tamen triginta annis non erit necesse inspicere hanc tabulam, sed solum indicem, qui post eam positus est. %sic
qui est a 1564, vsque ad 1593.




%pp28
\newpage
\color{red}
\begin{center} \hypertarget{TABULAFESTORUM}
 Tabula festorum mobilium iuxta vsum Roman\ae \ ecclesi\ae .
\end{center}
\vspace{-2em}
\bookmark[dest=TABULAFESTORUM]{TABULA FESTORUM MOBILIUM}

\color{black}
\begin{center}
\begin{tabular}{| l r | l | l | l | l | l | l | l | l |}
\hline
% \multicolumn{2}{l}{\color{red}Au. nu.} & & & \color{red} Dies & \\
\textfrak{Au.} & \textfrak{lta} & \textfrak{lxx.} & \textfrak{Linis} & \textfrak{pasch.} & \textfrak{ascensi.} & \textfrak{pente.} & \textfrak{cor. do.} & \textfrak{dominice} & \textfrak{aduentus}\\
\textgoth{Au.} & \textfrak{lta} & \textfrak{lxx.} & \textfrak{Linis} & \textfrak{pasch.} & \textfrak{ascensi.} & \textfrak{pente.} & \textfrak{cor. do.} & \textfrak{dominice} & \textfrak{aduentus}\\
\textswab{Au.} & \textfrak{lta} & \textfrak{lxx.} & \textfrak{Linis} & \textfrak{pasch.} & \textfrak{ascensi.} & \textfrak{pente.} & \textfrak{cor. do.} & \textfrak{dominice} & \textfrak{aduentus}\\



\end{tabular}

\end{center}









% \mainmatter


\newpage
\thispagestyle{empty}
\mbox{}
\newpage

\pagenumbering{arabic}

\chapter{Fundamentals of Vectors}

\section{Vectors in Euclidean Space}



\end{document}

